\documentclass[a4paper,fleqn]{cas-sc}

\usepackage[T5]{fontenc}
\usepackage[utf8]{inputenc}
\usepackage[vietnamese]{babel}
\usepackage[numbers]{natbib}
\usepackage{amsmath,amssymb}
\usepackage{graphicx}
\usepackage{booktabs}
\usepackage{url}

% Ghi chú người phụ trách (xóa lệnh này và các dòng \todo khi viết xong)
\newcommand{\todo}[1]{\par\noindent\textit{[#1]}\par}

\begin{document}

\let\WriteBookmarks\relax
\def\floatpagepagefraction{1}
\def\textpagefraction{.001}

\shorttitle{Antibandwidth Problem với SAT Solver và CP-SAT}
\shortauthors{Tường et al.}

\title[mode=title]{Giải bài toán Antibandwidth (ABP) bằng SAT Solver và CP-SAT}

\author{Nguyễn Xuân Tường}
\author{Hoàng Lâm Tùng}
\author{Doãn Duy Lợi}

\address{Khoa Công nghệ Thông tin, Trường Đại học Công nghệ, Đại học Quốc gia Hà Nội}

\begin{abstract}
\todo{Tường: Viết tóm tắt (Abstract) ở Tuần 3: Tổng quan bài toán ABP, tiếp cận mô hình hóa SAT (Direct, Order) và đối chuẩn CP-SAT, các kết quả thực nghiệm nổi bật trên tập Harwell--Boeing và kết luận.}
\vspace{6\baselineskip}
\end{abstract}

\begin{keywords}
Antibandwidth \sep
SAT solver \sep
Direct Encoding \sep
Order Encoding \sep
CP-SAT \sep
Symmetry Breaking \sep
cnsplots \sep
Harwell--Boeing
\end{keywords}

\maketitle

%% ------------------------------------------------------------

\section{Tổng quan bài toán Antibandwidth}
\todo{Tường: Phần mở đầu (Tuần 3) và Chương 1 (Tuần 1).}

\subsection{Giới thiệu}
\todo{Tường: Giới thiệu bài toán Antibandwidth (ABP), bối cảnh nghiên cứu và tầm quan trọng trong tối ưu hóa tổ hợp.}

\subsection{Định nghĩa hình thức bài toán ABP}
\todo{Tường: Song ánh $f: V \to \{1, \dots, n\}$, hàm mục tiêu $B_A(G) = \max_f \min_{\{u,v\} \in E} |f(u) - f(v)|$.}

\subsection{So sánh với bài toán Bandwidth}
\todo{Tường: Phân tích tính đối ngẫu giữa Bandwidth Problem (min-max) và Antibandwidth Problem (max-min).}

\subsection{Độ phức tạp tính toán (NP-hard)}
\todo{Tường: Phân tích độ phức tạp tính toán của ABP và chứng minh NP-hard trên đồ thị tổng quát.}

%% ------------------------------------------------------------
\section{Ứng dụng thực tế}
\todo{Tường: Chương 2 (Tuần 1).}

\subsection{Phân bổ tần số vô tuyến (Radio Frequency Assignment)}
\todo{Tường: Giãn cách tần số giữa các trạm phát sóng liền kề để giảm thiểu hiện tượng nhiễu giao thoa.}

\subsection{Giảm nhiễu xuyên âm trong thiết kế VLSI (Crosstalk Minimization)}
\todo{Tường: Tối ưu khoảng cách hình học giữa các đường mạch dẫn tín hiệu song song trong vi mạch tích hợp.}

\subsection{Lập lịch mạng cảm biến và kênh truyền}
\todo{Tường: Phân chia khe thời gian (time-slot) truyền thông nhằm tránh xung đột dữ liệu trong mạng không dây.}

\subsection{Bảo mật dữ liệu và định tuyến ẩn danh}
\todo{Tường: Phân tán vị trí lưu trữ và định tuyến mạng phân tán chống lại các tấn công phân tích luồng dữ liệu.}

%% ------------------------------------------------------------
\section{Chuyển đổi bài toán tối ưu sang bài toán quyết định}
\todo{Tường: Chương 3 (Tuần 2).}

\subsection{Bài toán quyết định $ABP(G,k)$}
\todo{Tường: Định nghĩa bài toán quyết định: kiểm tra sự tồn tại của song ánh $f$ sao cho $\forall \{u,v\} \in E, |f(u) - f(v)| \ge k$.}

\subsection{Cận trên và cận dưới của $k$}
\todo{Tường: Xác định chặn cận lý thuyết $k_{min} = 1$ và $k_{max} = \lfloor n/2 \rfloor$.}

\subsection{Thuật toán Binary Search}
\todo{Tường: Thiết kế và phân tích sơ đồ khối thuật toán tìm kiếm nhị phân giá trị $k$ tối ưu kết hợp SAT Solver.}

\subsection{Incremental SAT}
\todo{Tường: Tối ưu hóa việc giải gia tăng (Incremental SAT) sử dụng assumptions để tái sử dụng learned clauses qua các vòng lặp kiểm tra $k$.}

%% ------------------------------------------------------------
\section{Mô hình hóa bài toán (SAT \& CP-SAT)}
\todo{Tùng: 4.1, 4.2 và 4.5. Lợi: 4.3. Tường: 4.4 (Tuần 2).}

\subsection{Direct Encoding và hệ ràng buộc song ánh}
\todo{Tùng: Ma trận biến Boolean $x_{i,p}$, các ràng buộc ALO/AMO cho đỉnh và vị trí, tích hợp CardEnc (Sequential Counter / Commander qua PySAT).}

\subsection{Ràng buộc khoảng cách (Distance Clauses)}
\todo{Tùng: Hệ mệnh đề xung đột khoảng cách $\neg x_{u,p} \lor \neg x_{v,q}$ với mọi $\{u,v\} \in E$ và $|p - q| < k$.}

\subsection{Kỹ thuật phá vỡ đối xứng (Symmetry Breaking)}
\todo{Lợi: Ghim đỉnh có bậc lớn nhất $v_0$ vào nửa đầu dải vị trí $\neg x_{v_0, p}$ ($\forall p > \lceil n/2 \rceil$), chứng minh giảm 50\% không gian tìm kiếm.}

\subsection{Order Encoding}
\todo{Tường: Biến mệnh đề $y_{i,p} \iff f(i) \le p$, hệ ràng buộc đơn điệu và lan truyền khoảng cách theo mô hình Tamura et al.}

\subsection{Mô hình quy hoạch ràng buộc (Google OR-Tools CP-SAT)}
\todo{Tùng: Biến miền nguyên IntVar $f(u) \in \{1, \dots, n\}$, ràng buộc AddAllDifferent và ràng buộc khoảng cách $|f(u) - f(v)| \ge k$.}

%% ------------------------------------------------------------
\section{Kết quả thực nghiệm}
\todo{Lợi: Đặc tả dữ liệu, biểu đồ cnsplots, phân tích so sánh và đánh giá (Tuần 3). Tùng: Bảng số liệu thực nghiệm CP-SAT. Tường: Thu thập ma trận lớn.}

\subsection{Đặc tả dữ liệu Harwell--Boeing}
\todo{Lợi: Bảng thông số tập ma trận chuẩn Harwell--Boeing (số đỉnh $n$, số cạnh $m$, mật độ đồ thị, phân bố bậc đỉnh).}

\subsection{Môi trường và cấu hình thực nghiệm}
\todo{Lợi: Thông số phần cứng CPU, RAM, hệ điều hành, phiên bản Python, PySAT (Glucose4, CaDiCaL) và Google OR-Tools CP-SAT.}

\subsection{Kết quả và so sánh}
\todo{Lợi: Biểu đồ trực quan hóa bằng thư viện cnsplots (thời gian giải, số lượng clause, speedup) so sánh giữa Direct Encoding, Order Encoding và CP-SAT.}

\subsection{Phân tích và đánh giá}
\todo{Lợi: Phân tích ưu nhược điểm thực tế giữa SAT và CP-SAT theo mật độ đồ thị; đánh giá hiệu quả giảm thời gian của Symmetry Breaking và Order Encoding.}

\subsection{Hướng phát triển}
\todo{Lợi: Đề xuất mở rộng các ràng buộc toàn cục nâng cao, kỹ thuật phân rã đồ thị và áp dụng cho các biến thể bài toán dán nhãn băng thông.}

%% ------------------------------------------------------------
\section{Kết luận}
\todo{Lợi: Tổng kết các đóng góp chính của đề tài, đánh giá mức độ hoàn thành mục tiêu và kết luận chung (Tuần 3).}

%% ------------------------------------------------------------
\section*{CRediT authorship contribution statement}
\todo{Tường: Conceptualization, Methodology, Software (Graph Loader, Order Encoding), Writing - Original Draft (Chương 1, 2, 3, 4.4). Tùng: Software (Direct Encoding, CP-SAT), Writing - Original Draft (Chương 4.1, 4.2, 4.5). Lợi: Software (Solver Pipeline, Symmetry Breaking, cnsplots Visualizer), Writing - Review \& Editing (Chương 4.3, 5, 6).}

\section*{Declaration of competing interest}
\todo{Tường, Tùng, Lợi: The authors declare that they have no known competing financial interests or personal relationships that could have appeared to influence the work reported in this paper.}

\section*{Data availability}
\todo{Tường: Toàn bộ mã nguồn, dữ liệu thực nghiệm và tài liệu hướng dẫn được công khai trên GitHub tại \url{https://github.com/LoiDoancode/Antibandwidth_Problem} (Branch: \texttt{Loi\_ABP}).}

%% ------------------------------------------------------------
\appendix
\section{Phụ lục}
\todo{Tường, Tùng, Lợi: Bảng số liệu chi tiết toàn bộ các ca kiểm thử benchmark, mã lệnh tái lập kết quả và hướng dẫn sử dụng CLI.}

%% ------------------------------------------------------------
\nocite{*}
\begin{thebibliography}{1}
\bibitem{tamura2009}
\todo{Tường: Trích dẫn bài báo Order Encoding (Tamura et al., 2009).}
\bibitem{harwellboeing}
\todo{Tường: Trích dẫn tập dữ liệu chuẩn Harwell--Boeing Matrix Collection.}
\bibitem{ortools}
\todo{Tường: Trích dẫn Google OR-Tools CP-SAT Solver.}
\bibitem{pysat}
\todo{Tường: Trích dẫn thư viện PySAT.}
\bibitem{cnsplots}
\todo{Tường: Trích dẫn thư viện trực quan hóa cnsplots.}
\end{thebibliography}

\end{document}

